% Lösungen Einheit 5
\einheitenkopf*{Einheit 5 von 5 · Potenzfunktionen mit natürlichem Exponenten}

\erg{43}{a) $0;\ 1;\ 8;\ 27$ \quad b) $-27;\ -8;\ -1$ \quad c) $16;\ 1;\ 0;\ 1;\ 16$ \quad d) $0{,}0625;\ 0{,}0625;\ 81$}
\erg{44}{a) $8$ \quad b) $16$ \quad c) $27$ \quad d) $10\,000$ \quad e) $-125$ \quad f) $16$ \quad g) $-16$ \quad h) $0{,}125$ \quad i) ja, $(-2)^4 = 16$ \quad j) nein, $3^3 = 27 \neq -27$ \quad k) $x = 6$ \quad l) $x = 3$ oder $x = -3$}
\erg{45}{a) Parabel durch $(-2\,|\,4)$, $(-1\,|\,1)$, $(0\,|\,0)$, $(1\,|\,1)$, $(2\,|\,4)$ \quad b) durch $(-2\,|\,{-8})$, $(-1\,|\,{-1})$, $(0\,|\,0)$, $(1\,|\,1)$, $(2\,|\,8)$ \quad c) durch $(-2\,|\,16)$, $(-1\,|\,1)$, $(0\,|\,0)$, $(1\,|\,1)$, $(2\,|\,16)$ – flacher um $0$, steiler außen als die Parabel}

\begin{ksys}[xmin=-2,xmax=2,ymin=-8,ymax=16,xstep=0.5,ystep=2,klein]
\funktion[1.8]{\x^2}{x^2}
\funktion[1.6]{\x^3}{x^3}
\funktion[-1.7]{\x^4}{x^4}
\end{ksys}

\erg{46}{a) zur $y$-Achse \quad b) zum Ursprung \quad c) zur $y$-Achse \quad d) zum Ursprung \quad e) schmaler \quad f) breiter \quad g) an der $x$-Achse gespiegelt, der Graph ist nach unten geöffnet}
\erg{47}{a) $y = x^3$: II \quad b) $y = -x^2$: III \quad c) $y = x^4$: I \quad d) $y = 0{,}5x^3$: IV \quad e) I (gerader Exponent, positive Werte, symmetrisch zur $y$-Achse)}
\erg{48}{a) Potenzfunktion \quad b) Exponentialfunktion \quad c) Potenzfunktion \quad d) Exponentialfunktion \quad e) Exponentialfunktion \quad f) Potenzfunktion \quad g) $x = 2$: $8$ und $9$; $x = 5$: $125$ und $243$ – $3^x$ wächst schneller}
\erg{49}{a) Fehler: Tom rechnet $3 \cdot x$ statt $x \cdot x \cdot x$. Richtig: $f(-2) = (-2)^3 = -8$ \quad b) $-27$ \quad c) Fehler: Sara setzt $-2$ ein. $-x^4$ heißt $-(x^4)$: $f(2) = -(2^4) = -16$ \quad d) $-81$}
\erg{50}{a) $x^4 = x \cdot x \cdot x \cdot x$ hat vier gleiche Faktoren: Bei negativem $x$ heben sich je zwei Minuszeichen auf, $0^4 = 0$; das Ergebnis ist nie negativ. \quad b) Für negative $x$ ist $x^3$ ein Produkt aus drei negativen Zahlen und damit negativ.}
\erg{51}{a) $7^3 = 343\,\text{cm}^3$ \quad b) $x^3 = 1\,331$, also $x = 11\,\text{cm}$ \quad c) $(2x)^3 = 8x^3$: Das Volumen wird achtmal so groß.}
